% GENERATED, edit the French unit or the translation table
% GENERATED from content/data/enrichment.json by mxbuild.enrichment; do not edit.
\mxunit{enrichment-telescoping-notes}{Telescoping sums}
\mxsubtitle{High school · sequences \textperiodcentered{} Lesson and method \textperiodcentered{} 10}

\begin{prerequisites}
Fractions, finite sums and indices.
\end{prerequisites}

\begin{keypoint}[Ideas and method]
A telescoping sum is a sum of consecutive differences. Interior terms cancel: $\sum (a_{k}-a_{k+1})=a_{1}-a_{n+1}$ for k from 1 to n.

Write the first three terms and the last two to expose cancellation and check endpoints. Term-by-term simplification does not remove the need to specify indices.

A useful decomposition is $\frac{1}{k(k+1)}=\frac{1}{k}-\frac{1}{k+1}$, valid for $k\geq 1$. Verify it by taking a common denominator.
\end{keypoint}

\begin{example}
$\sum  \frac{1}{k(k+1)}$, from $k=1$ to n, equals $1-\frac{1}{n+1}=\frac{n}{n+1}$. The sum is below 1 and tends to 1 as n grows.
\end{example}

\begin{activity}[Explore]
Compute the first five partial sums, then study the gap to 1. Distinguish the finite sum from its limit.
\end{activity}

